% !TEX root = ./main.tex

\chapter{绪论}

\section{研究背景与意义}

近年来，生成式人工智能技术快速发展，以扩散模型、生成对抗网络和多模态大模型为代表的AIGC（AI Generated Content）技术在图像生成、视频生成、虚拟人驱动、内容编辑等任务中表现出较强的内容生成能力。AIGC技术一方面提高了数字内容生产效率，推动了媒体制作、教育培训、广告营销和智能交互等领域的发展；另一方面，也带来了虚假信息传播、深度伪造、不安全内容生成、身份冒用以及平台审核压力增加等一系列安全问题。特别是在图像与视频场景中，AI生成内容已经逐步具备较强的逼真性和传播性，使得传统依赖人工经验的审核方式面临准确率低、效率不足和主观性较强等局限。

在此背景下，围绕AIGC内容的真实性与安全性开展研究具有重要意义。从理论层面看，系统梳理图像与视频生成内容中的安全风险类型，分析不同检测方法的适用范围与技术特点，有助于丰富AIGC安全治理的研究体系；从应用层面看，构建面向真实场景的AI内容安全检测与分析工具，有助于为网络内容审核、媒体真实性核验、平台合规治理以及个人权益保护提供技术支持。因此，设计并实现一个兼顾虚假内容检测与不安全内容识别的多模态安全分析平台，具有较强的研究价值和现实意义。

\section{国内外研究现状}

\subsection{AIGC生成技术研究现状}

AIGC技术的发展为图像和视频内容生成提供了新的技术基础。图像生成方面，基于扩散模型的Stable Diffusion等方法在生成质量、细节还原和文本对齐能力方面取得了显著进展，推动了文生图、图像编辑和风格迁移等任务的广泛应用。视频生成方面，文本到视频、图像到视频和数字人动画等方向也不断发展，典型方法在时序一致性、运动建模和多模态对齐方面持续提升。与此同时，Deepfake、人脸替换和属性编辑等技术的普及，使得伪造内容的构造门槛明显降低，进一步加剧了内容真实性识别的难度。

\subsection{AI生成内容检测研究现状}

针对AI生成图像和视频的检测问题，国内外研究者提出了多种技术路线。早期方法主要基于图像纹理、频域特征和生成伪影进行判别，这类方法在特定生成模型下具有较好效果，但在跨模型泛化方面存在不足。随着深度学习的发展，研究者提出了基于CNN的分类检测方法，通过学习真实样本与伪造样本之间的差异实现自动判别。进一步地，利用大规模预训练视觉模型特征的通用检测方法逐渐受到关注，例如UniversalFakeDetect利用预训练视觉表征与轻量分类头相结合，在跨生成模型检测中表现出较好的泛化能力。除此之外，针对扩散模型图像的重建型、对比学习型检测方法，以及面向视频深度伪造的训练自由方法和局部全局联合分析方法，也成为近年来的重要研究方向。

总体来看，现有检测方法在特定任务上取得了一定成效，但仍面临若干问题：其一，不同方法的适用对象存在差异，部分方法主要面向图像，难以直接推广至视频场景；其二，本地模型虽然具备自主可控和便于实验复现的优点，但部署复杂度较高，对算力资源和工程实现要求较强；其三，多模态大模型API具备较好的开放性和工程可用性，但往往依赖外部平台，成本、稳定性和可解释性仍有待进一步平衡。因此，如何结合本地模型与外部能力，构建统一、可扩展的检测框架，是当前AIGC安全研究与工程应用中的一个重要问题。

\subsection{不安全内容审核研究现状}

除真实性检测外，AIGC内容还可能涉及暴力、色情、仇恨、敏感符号等不安全信息，这些问题同样是当前内容治理的重点。传统的不安全内容识别多依赖图像分类、目标检测或多标签识别模型，对固定类别场景具有较好性能。近年来，多模态大模型逐步被引入内容审核任务，能够结合视觉信息与语义提示完成更灵活的风险识别与结果解释。此外，一些训练自由的安全防护方法也为生成内容的前置治理提供了新的思路。然而，从实际应用角度看，单一模型往往难以兼顾识别范围、工程效率与部署成本，因此构建多源能力融合的安全分析平台具有较高的现实价值。

\section{本文的主要研究内容}

围绕AIGC图像与视频内容中的虚假风险和不安全风险，本文以“检测方法整合与平台化落地”为主线，设计并实现了一个面向多模态场景的AI内容安全检测与分析工具。本文的主要研究内容包括以下几个方面。

第一，围绕AIGC安全风险开展调研与分析。本文对AI生成图像、视频及深度伪造内容中的主要风险类型进行了梳理，重点分析了虚假生成、不安全内容和平台审核场景下的检测需求，总结了现有方法在适用范围、部署方式和工程复杂度等方面的特点。

第二，整合多源检测能力，构建统一检测框架。针对图像虚假检测任务，本文接入了本地部署的UniversalFakeDetect模型，并结合多模态大模型API实现图像与视频的AI生成鉴别；针对不安全内容识别任务，本文接入图像和视频审核API，实现对违规内容的自动化检测。通过统一的后端代理接口和结果解析逻辑，平台能够对不同来源的检测能力进行封装与调度。

第三，设计并实现AI内容安全检测与分析平台。本文采用前后端分离的系统架构，前端使用React与TypeScript实现交互界面和结果展示，后端使用Node.js与Express实现接口代理、鉴权签名、任务轮询和业务编排，本地模型服务采用Python与FastAPI实现推理封装。系统支持图片与视频输入、检测结果可视化、风险分级展示以及内容管理等功能，形成了较为完整的工程闭环。

第四，结合实验样本对系统功能与检测效果进行验证。本文通过本地生成模型与外部生成能力构造测试样本，并结合公开或真实样本对检测模块进行功能验证和结果分析，从而评估不同检测方式在工程应用中的适用性与局限性。

\section{本文的特色与创新点}

相较于单一模型复现或单一能力调用，本文的工作特色主要体现在以下几个方面。

第一，面向图像与视频场景，构建了兼顾虚假内容检测与不安全内容识别的多模态安全分析平台。该平台不仅关注AI生成内容是否为伪造，还关注其是否包含违规与不当信息，从而更贴近真实应用场景下的内容治理需求。

第二，提出并实现了本地模型与外部API相结合的统一检测框架。通过后端代理层对鉴权、转发、异步任务轮询和结果格式进行统一封装，平台能够兼容多类检测能力，在保证工程可用性的同时提升系统扩展性。

第三，实现了从内容生成、内容检测到结果展示与管理的闭环式工具平台。该设计既能支持实验样本构建与方法验证，也便于后续扩展更多检测模型与治理能力，具有较好的工程实践价值。

需要说明的是，本文的主要贡献集中在系统设计、方法整合、工程实现与实验验证方面，而非提出新的底层检测算法。

\section{论文结构安排}

本文共分为六章，各章内容安排如下：

第一章为绪论，介绍研究背景、研究意义、国内外研究现状、本文的主要研究内容、特色与创新点，以及论文整体结构安排。

第二章为相关技术与方法综述，主要介绍AIGC生成技术、AI生成内容检测方法、不安全内容审核方法以及系统实现所涉及的关键技术，为后续系统设计与实现奠定理论基础。

第三章为系统需求分析与总体设计，结合实际应用需求对系统功能与非功能需求进行分析，给出平台的总体架构、模块划分和核心数据流设计。

第四章为系统详细设计与实现，重点介绍前端交互界面、后端代理服务、本地模型服务、外部API接入方式以及检测结果可视化与系统部署实现。

第五章为实验设计与结果分析，介绍实验环境、测试样本、评价指标以及虚假内容检测、不安全内容检测和系统性能方面的实验结果，并对结果进行分析讨论。

第六章为总结与展望，对全文工作进行总结，分析当前系统仍存在的不足，并对后续研究方向进行展望。
